\documentclass[11pt]{article}
\usepackage{../shared/luxcover}
\usepackage{amsmath,amssymb}
\usepackage{hyperref}
\usepackage{booktabs}
\usepackage{enumitem}
\usepackage{xcolor}

\title{Lux Governance}
\author{Zach Kelling\\\textit{Lux Industries, Inc.}\\\texttt{research@lux.network}}
\date{April 2024}

\begin{document}

\luxcoverpage

\begin{abstract}
Decentralized Governance for LUX
\end{abstract}

\tableofcontents
\newpage


\section{Abstract}

Corporate entities of all types are governed by rules describing permissible and prohibited behavior. These regulations may exist as private contracts (such as bylaws or shareholder agreements) between business shareholders. They may also be imposed by law in addition to or in lieu of a formal agreement between the parties.


Historically, corporations have only been able to operate through people or through corporations that are ultimately owned by people. This raises two elementary and fundamental issues. Whatever a private contract or public law may require, (1) people do not always follow the rules, and (2) people do not always agree on what the rules require. Collaboration without a corporate structure does not necessarily alleviate existing issues, and may even introduce new ones. In the absence of a corporate form, an explicit written agreement replaces ambiguous informal "understandings," and the legal safeguards afforded by a corporate form are unavailable.


Rule-breaking within an organization is not always visible, and shareholder motivations may be irrelevant once their money has been lost. Although a firm or its management may be held civilly or criminally accountable for bad behavior, investors who have already lost money may find the punishment to be of little consolation.


The situation is exacerbated by crowdfunding. On the one hand, it has made it simpler for small contributors to invest in huge initiatives, and on the other hand, it has made it feasible for entrepreneurs to acquire financial backing that before may not have been readily accessible. On the other hand, small investors remain susceptible to financial mismanagement or blatant fraud, and because they have a limited ownership in a company, they may lack the ability to recognize concerns, participate in governance choices, or recover their investment with relative ease.


Similarly, corporate leadership and management may be accused of malfeasance or mismanagement even when they think they have behaved in good faith based on their understanding of their obligations and interpretation of applicable rules.


The Lux DAO infrastructure provides a Governance solution based on Lux Network, a blockchain platform that combines a Turing-complete programming language with smart contract processing functionality coupled with a money transmitters license and a zero knowledge privacy bridge.


This paper describes a mechanism that, for the first time, leverages holographic consensus, quadratic voting, aged staking and weighted POS based on experience points gained through reputation and verified with proof of human biometric digital signal processing to authenticate user identity.


This enables the formation of organizations in which participants keep direct, real-time control over contributed cash and governance rules are written, mechanized, and enforced using software. Specifically, standard smart contract code that can be used to create a Decentralized Autonomous Organization (DAO) that can be published on the Lux blockchain. This paper illustrates how the Lux DAO's code operates, with a focus on several fundamental formation and governance characteristics, such as structure, creation, and voting rights.


\section{Introducing Lux Triforce Voting}

Lux DAO is the portal that enables access to Lux Network's governance modules. The governance modules comprise the ability to vote, ability to stake your LUX, as well as the ability to participate in open governance. The Lux DAO is an important part of the Lux ecosystem because it helps to keep the network decentralized and allows users to have a direct impact on the direction of the project.


The Lux DAO is a decentralized autonomous organization that is responsible for the governance of the Lux Protocol. The Lux Protocol is a set of smart contracts that are used to create and manage tokenized assets. The DAO makes decisions about the Lux Protocol through voting by its members.


By joining the Lux DAO, you can help make decisions about the Lux Protocol and earn rewards in the form of KARMA. You can also participate in the Lux Protocol by creating and managing tokenized assets. If you are interested in becoming a member of the Lux DAO, acquire one of our NFTs to unlock access.


The pairing  of LUX  and KARMA are combined with a Stake Age duration to form the core of the Lux Governance and Voting Protocol, FLUX. These three factors will determine your DAO Rank, which in turn determines how many FLUX tokens you receive to cast votes as well as fee rewards you earn for staking your tokens.


\section{Benefits of Staking}

By staking KARMA, holders can advance in The Lux DAO's rankings. Achieving a high rank offers numerous advantages, including the ability to cast more votes, collect more fee rewards, and get dApp-level privileges. Each rank in the DAO is associated by a level of Network responsibility. Users are held accountable for their activities with the possibility of losing their stake position, its value, and the rights connected with it. To summarize, this means staking can be seen as a person's proof of commitment and it is used as a primary quality measure in the network.


\subsection{The Lux DAO voting process is done using the following mechanisms;}

Lux decentralized Identification service: Each community member must undergo a verification process before voting on Dao proposals. This is done through the Lux decentralized Identification service.


Holographic Consensus: the DAO's Creation Phase and fundamental operations are presented through holographic consensus.


Experience Weighted Voting: The cost of votes are weighted against experience points that community members earn as they participate in community discussions and social platform interactions earning them reputation as they continue to make voting predictions that result in positive outcomes for the community.


Quadratic Voting:  after a proposal is approved by the Dao committee, verified  Lux DAO community members are given the ability to purchase votes for each proposal based on a quadratic squared value equation.


\section{Lux Decentralized Identification Services}

There are a number of ways to cheat the Lux Dao voting system when it comes to transferring private keys to sell off voting rights and Quadratic voting makes token-selling even more lucrative since first votes cost significantly less then every vote purchased thereafter. The Lux DAO mitigates this issue by employing identification procedures that make it difficult to transfer voting rights. Prior to voting, users must register a profile with the Lux decentralized Identification Services. The Lux DIS is a combination of NFC (near field communication coupled with a DSP (digital signal processing) machine learning algorithm to authenticate users' Identification and even in some cases run it against background checks in compliance with KYC and AML regulatory requirements.


\section{Holographic Consensus}

Holographic consensus blends features of prediction markets with DAO governance.


Holographic consensus is a voting system used by Lux DAO to help scale governance. Currently, there are DAOs with thousands of members, making it difficult, if not impossible, to get members to vote on every proposal.


Holographic consensus is rational for multiple reasons. It decreases voter fatigue, which arises when people are obliged to vote too frequently. In governance provided by HC, DAO members can focus on proposals with high stakes and make decisions swiftly.


The rationale behind holographic consensus is that attention spans are limited, thus only the most crucial propositions should be put to a vote. Unnecessary proposals can be sorted out, ensuring that DAO members' attention is properly focused.


\subsubsection{How It Works}

Governance is made possible through Lux DAO's Governance token, KARMA. Token holders cannot vote on ideas, but they can "bet" on which proposals have the best chance of success.


If the proposal is approved, holders will receive more KARMA tokens. But if the initiative fails, their KARMA holdings will be reduced.


Any DAO member who meets the necessary requirements may submit a proposal. Proposals that receive a large number of KARMA bets are "boosted" and advance to the next round. At this stage, suggestions with modest KARMA wagers are eliminated.


Members with voting rights vote on Round 1 proposals. KARMA holders can either make a profit or incur a loss on their wagers, depending on the outcome of their choice proposal. Following approval, the DAO executes and implements the authorized ideas.


In addition, it stops hostile users from attacking a DAO by submitting malicious proposals. A potential attacker must fund the promotion of a proposal by purchasing KARMA tokens. This increases the cost of an attack and may deter governance-based DAO attacks.


\section{Voting Based on Reputation and Weighting}

\subsection{Why Use Reputation and Weighted Voting?}

\subsubsection{Rational Indifference}

If tokens held = voting power, then DAO members with less tokens may have fewer incentives to participate in governance. A person who believes their influence on the result of a proposition to be negligible may rationalize their unwillingness to vote.


The phenomenon known as "rational indifference" can deter DAO members from active involvement, putting token whales in charge of governance. This reintroduces the principal-agent problem and undermines the core value proposition of DAOs, which is to place power in the hands of all stakeholders.


This creates an issue that allows crypto whales and other outsiders purchasing governance tokens to exert disproportionate influence in a DAO. In this situation, one individual uninterested in the long-term survival of the group can amass sufficient influence to destabilize it for personal gain.


The Lux DAO employs weighted and reputation-based voting techniques to prevent this issue. These two notions share some similarities, but they are not identical.


\subsubsection{Weighted Voting Systems}

In weighted voting systems, the voting power of tokens is proportional to the length of time they have been kept. Weighted voting systems motivate token holders to lock up their tokens in order to increase the value of their votes.


In a DAO, weighted voting techniques make it more difficult to acquire voting power. You cannot just borrow governance tokens from the market and use them to quickly affect vote outcomes. This will not completely eliminate governance assaults, but it will considerably increase the expense of performing one.


\subsubsection{Reputation}

The voting mechanism based on reputation follows a similar trend. Members of the DAO receive a single, non-transferable token for voting on proposals. And the voting power associated with each token is tied to the holder's reputation.


Reputation points can be acquired in various ways. A DAO may award reputation points to members based on their membership history. Or, members can gain extra reputation points by positively contributing to the expansion of the organization (for example, by completing tasks as well as participating in community discussions and receiving upvotes for comments on proposals).


Voting based on a candidate's reputation helps ensure that individuals with a lengthy history of involvement in an organization attain positions of prominence. This indicates that it will be difficult for an outsider to enter the Lux DAO through the community governance processes.


\subsection{Unfair Representation}

In rare cases, reputation voting may be judged unfair. What if, for instance, a respectable member suddenly becomes evil and utilizes their influence inside the group for nefarious purposes? Wouldn't reputation systems merely revive the tragedy of the commons?


The fact that someone earned reputation points in the early stages of a project does not imply that they should always have leadership status. It would be unfair for genesis members to have greater authority than new, active members, especially if their contribution decreases.


\subsection{Slashing and Expiration}

Lux proposes two solutions to this issue. First, DAOs may permit the revocation of reputation points. If a member commits wrongdoing, their reputation points should decrease accordingly.


Second, it is possible to program reputation points to decrease over time or even expire. This solves the problem of injustice and maintains the fluidity of power in DAOs.


However, conventional reputation systems have additional drawbacks such as Anyone can transfer their reputation by selling their private keys (reputation points are linked to wallets). Despite the difficulty of establishing a healthy market for the selling of reputation points, this mechanism is viable for supporting competent DAO governance. To mitigate against this type of attack on the system, Lux Dao requires each voting participant to verify the authentication of their identity through the Lux Decentralized Identification service.


\section{Quadratic Voting}

Quadratic voting is an innovative governance system that has the endorsement of Ethereum's creator Vitalik Buterin. Quadratic voting is attractive because it provides an alternative to token holder voting and lessens the impact of whales on voting outcomes.


In a quadratic voting system, each individual has one vote on a proposition. However, users may purchase additional "credits" in order to cast further votes for their favorite alternative.


When comparing quadratic voting systems to conventional POS where the cost of additional votes is linear in a 1T1V (one token/ one vote) system, quadratic systems are much different. In a one token / one vote system, If you need five additional votes, you simply need five additional tokens. The cost of extra votes in a quadratic voting system is quadratic; one vote costs one token, two votes cost four tokens, three votes cost nine tokens, four votes cost sixteen tokens and so on and so forth.


Individuals can express their level of support for a proposal by purchasing vote credits through quadratic voting. This makes quadratic voting the equivalent of "putting your money where your mouth is" within the DAO.


Concurrently, it decreases the impact of crypto whales on decision-making. It is evident from our earlier illustration that the usefulness of more voting diminishes while the costs increase. A bad actor within a DAO would have to spend a lot of money to obtain enough tokens to manipulate outcomes.


\section{Dao Ranking}

Your DAO Rank will be useful for, among other things, proposal writing, vote allocations, fee distributions, etc. The precise DAO-level specifications will be determined in the technical paper. By using an exponential function to determine the minimum requirement of KARMA Power, it will be easier for newcomers with lower rankings to participate in the DAO. It will, however, take longer if you want to participate with higher levels of DAO ranking.


\section{LUX Power}

LUX Power is a factor that determines a participants weight and rank within the DAO. The Lux Staking Protocol produced this measurement. You can enhance your LUX Power in two ways:


Increase the number of staked LUX and KARMA tokens


Stake your tokens longer to increase their Stake Age


The first is very straightforward: the more tokens you stake, the greater your KARMA Power. This is a straightforward 1:1 ratio (e.g., 100k KARMA staked = 100k KARMA Power). This is how the majority of common staking protocols operate. However, The Effect Staking Protocol provides a measuring option to track the duration of your active stake.


This attribute is known as Stake Age.


\section{Stake Age}

Stake Age is a factor that gradually increases the weight of staked tokens over time; the longer tokens are staked, the more KARMA Power they will possess.


Stake Age served as a multiplier for the number of FLUX tokens generated by pairing LUX and KARMA tokens. The longer the Stake Age, the greater the amount of FLUX that will be generated in that time.


Stake Age has various benefits above normal staking:


Users will be incentivized to stake their tokens for longer periods of time.


Early and sustained participation is rewarded.


New members with a large number of tokens will have a more difficult time instantly claiming large rewards or votes from the DAO.


Users will lose their Stake Age for unstaked tokens, increasing the penalty for un-staking and incentivizing users to keep their tokens staked.


Voting Allotments and Fee Incentives


\section{Voting Allocations and Fee Rewards}

There are four major stages that play an important role in the Lux DAO. The first one is the Staking Protocol, which will determine your DAO Rank. The next stage is to become a DAO Guardian with your DAO Rank by signing the DAO Constitution and participating in the Network. As a DAO Guardian, you can use your DAO Rank to write proposals and cast votes. By participating in the DAO, you'll earn stakes in the network fee distribution pool which is the last stage of the Lux DAO.


\subsection{The four stages of the Lux DAO}

The number of votes that a Guardian has is determined by the Lux Staking Protocol: a Triforce of LUX and KARMA along with the choice of a minimum one year Stake Age up to a seven year Stake Age. These three factors determine your DAO Rank and how many voting allocations you have at your disposal. Participating in either voting or proposing is highly rewarded with fee distributions, which will be discussed in a forthcoming chapter. Tied Voting Allocations will also play a special role in the fee distribution later on.


While this paper outlines the high-level inner workings of the Lux Staking Protocol, the exact implementation, parameters, and details are being finalized in a Technical Paper.


The technical paper will cover the mathematical formulas, among other things, for KARMA Power with limited Stake Age, the DAO level requirement formulas, dilution of the Stake Age when topping up your stake, voting allocations, and more. Keep an eye out for this Technical Paper if you are interested in learning more about the specifics.


\end{document}
